\documentclass[10pt,a4paper]{amsart}
\usepackage[margin=19mm]{geometry}
\usepackage{amsmath,amssymb,mathtools}
\usepackage[scr=rsfs,cal=euler]{mathalfa}
\usepackage{xfrac}
\usepackage[unicode,hidelinks,bookmarks=false]{hyperref}
\newcommand{\ie}{i.e.\ }
\newcommand{\cf}{cf.\ }
\newcommand{\resp}{respectively\ }
\newcommand{\Subsection}{\subsection}
\providecommand{\nobreakdash}{\nobreak\hbox{-}}
\setlength{\emergencystretch}{2em}
\multlinegap0pt
\usepackage{bm}
\usepackage{bbm}
\usepackage{dsfont}
\usepackage{xspace}
\usepackage{cancel}
\usepackage{mathtools}
\usepackage{amsthm}
\makeatletter
\@ifundefined{theorem}{\newtheorem{theorem}{Theorem}}{}
\@ifundefined{lemma}{\newtheorem{lemma}[theorem]{Lemma}}{}
\makeatother
\DeclareMathOperator{\re}{Re}
\title[A Complex Analogue of Spencer's Six Standard Deviations Theo]{A Complex Analogue of Spencer's Six Standard Deviations Theorem and the Complex Banach-Mazur Distance}
\author{Tomasz Kobos, Marin Varivoda}
\date{Source: arXiv:2602.12868v1 (CC BY 4.0)}
\begin{document}
\maketitle

\noindent\emph{Source and licence.} A Complex Analogue of Spencer's Six Standard Deviations Theorem and the Complex Banach-Mazur Distance. Version arXiv:2602.12868v1 (CC BY 4.0), distributed under the Creative Commons Attribution 4.0 International licence, as recorded in the arXiv licence metadata for this exact version and shown on the abstract page https://arxiv.org/abs/2602.12868v1. Full rights notice: licenses/arxiv-2602.12868-CC-BY-4.0.txt.\par\smallskip

\noindent\textit{Editorial scope.} Counted unit: the Lemma whose source label is \texttt{lem:grid\_placement} in the article (source0.tex lines 447--449, source label \texttt{lem:grid\_placement}), delivered together with the article's own complete original proof of it (source0.tex lines 451--503, one \texttt{proof} environment of 53 lines). No mathematics is written, repaired or re-derived: statement and proof are the article's own text line for line. Required context retained uncounted: lem:blob\_expansion (lines 423--425). Every definition, display and label of those lines is retained unchanged. Omitted from this package and not redistributed: 1--422, 426--446, 450--450, 504--817. Nothing inside a retained excerpt is omitted or altered: every retained body is delivered exactly as the article's lines are written, so no omission receipt of any kind accompanies this package. Citations: the retained excerpts carry no citation command at all, so no bibliography entry of the article is transcribed and no reference is redistributed. Reference adaptation: none was needed and none was made; every reference inside the delivered unit resolves inside it, so no unresolved reference or citation is produced. Printed slips kept literal: no printed word, symbol, hypothesis or number of the counted statement or of its proof was altered. Layout and class: the article's own class is \texttt{article} with packages including inputenc, fontenc, subcaption, geometry, listings, xcolor, amsmath, amssymb, amsthm, graphicx, hyperref; the wrapper is the lane's pinned \texttt{amsart} editorial wrapper at 10pt on A4, which restates the theorem-like environments the retained text needs and adds the article's own macro definitions that the retained text needs (\texttt{\textbackslash re}), with extra wrapper packages (bm, bbm, dsfont, xspace, cancel, mathtools). The delivered numbering is the unit's own. Assets: no retained span contains a figure environment, a graphics-inclusion command, a PDF-page inclusion, a table or a TikZ picture; no image file is copied into this package, the package declares no asset dependency and no image file is redistributed with it. Licence: version arXiv:2602.12868v1 of this article is distributed under the Creative Commons Attribution 4.0 International licence, as recorded in the arXiv licence metadata for this exact version and shown on the abstract page https://arxiv.org/abs/2602.12868v1. A full rights notice is delivered at licenses/arxiv-2602.12868-CC-BY-4.0.txt. Characters: every retained excerpt is pure ASCII, so the delivered engine, XeLaTeX with its default Latin Modern text font, reports no missing character.
\begin{lemma}\label{lem:blob_expansion}
Let $a \in \mathbb{C}^3$ be a non-zero vector such that $\|a\|_{\infty} \leq 1$. Then, there exists a vector $v \in \mathbb{C}^3$ with at most one coordinate of modulus less than $1$ and such that $B(a) \subseteq B(v)$.
\end{lemma}

\begin{lemma} \label{lem:grid_placement}
For any nonzero vector $a=(a_1, a_2, a_3) \in \mathbb{C}^3$ with $\|a\|_{\infty} \leq 1$, there exists a unimodular vector $x=(x_1, x_2, x_3) \in \mathbb{C}^3$, such that for $v=(a_1x_1, a_2x_2, a_3x_3) \in \mathbb{C}^3$ the toric body $B(v)$ contains at most one point of the grid $\mathcal{G}$.
\end{lemma}

\begin{proof}


By Lemma \ref{lem:blob_expansion} we can assume that at most one coordinate of $a$ has modulus strictly less than $1$. Without loss of generality, let $|a_1|=b \leq |a_2|=|a_3|=1$. We shall prove that the desired vector $v$ can be chosen as $v=(b, u, \overline{u})$, where $u=\exp{\left (i\left ( \arccos{\frac{b}{2}} - \frac{\pi}{3} \right ) \right )}$. In this case, we will show that $(0, 0)$ is the only point of $\mathcal{G}$ that could possibly belong to $B(v)$.

For fixed $k, l \in \{0, 1, 2\}$ we have
$$\left | \left \langle (b, u, \overline{u}), (1, \omega^k, \omega^l) \right \rangle \right |^2 = b^2 + 2 + 2 \re\left ( b u \omega^{-k} + b u \omega^{l} + u^2 \omega^{-k} \omega^{l} \right ).$$
Thus, by changing $-k$ to $k$, what we need to prove can be simply restated as
\begin{equation}
\label{desiredineq}
b^2 + 2b \re \left (u \left ( \omega^{k} + \omega^{l} \right ) \right ) + 2 \re \left (u^2\omega^k\omega^l \right ) \leq 1
\end{equation}
for $k, l \in \{0, 1, 2\}$ and $(k, l) \neq (0, 0)$. To this end, we shall consider some different cases.

\begin{enumerate}
\item $(k, l)=(1, 2)$ or $(k, l) = (2, 1)$. Then $\omega^k + \omega^l=-1$ and $\omega^k \omega^l = 1$, so the desired inequality \eqref{desiredineq} takes the form
$$b^2 - 2b\re(u) + 2\re(u^2) \leq 1.$$
By using the definition of $u$, the fact that $b \in [0, 1]$ and some standard trigonometric formulas, we obtain that
$$b^2 - 2b\re(u) + 2\re(u^2) = b^2 - 2b \cos \left ( \arccos{\frac{b}{2}} - \frac{\pi}{3} \right ) + 2 \cos \left (2\arccos{\frac{b}{2}} - \frac{2\pi}{3} \right ) $$
$$=b^2 - 2b \left( \frac{b}{4} + \frac{\sqrt{3}}{2} \sqrt{1 - \frac{b^2}{4}} \right) + 2 \left( -\frac{1}{2} \left( \frac{b^2}{2} - 1 \right) + \frac{\sqrt{3}}{2} \sqrt{b^2 - \frac{b^4}{4}} \right)$$
$$=b^2 - \frac{b^2}{2} - b\sqrt{3\left ( 1 - \frac{b^2}{4} \right )} + \left ( 1 -\frac{b^2}{2} \right ) + \sqrt{3\left ( b^2 - \frac{b^4}{4} \right )} = 1.$$
So, in this case, the expression turns out to be actually equal to $1$ for all $b \in [0, 1]$.
\item $(k, l) = (1, 0)$ or $(k, l) = (0, 1)$. Then $\omega ^k + \omega^l = 1 + \omega = -\omega^2$ and $\omega^k \omega^l = \omega$, so the inequality \eqref{desiredineq} takes the form
$$b^2  - 2b \re \left (u \omega^2 \right ) + 2 \re \left (u^2\omega \right ) \leq 1.$$
Reasoning similarly as in the previous case, we get
$$b^2  - 2b \cos \left (\arccos{\frac{b}{2}} + \pi \right ) + 2 \cos \left (2\arccos{\frac{b}{2}} \right ) = b^2  + 2b \left ( \frac{b}{2} \right ) + 2 \left ( \frac{b^2}{2} - 1 \right ) = 3b^2 - 2,$$
which is at most $1$, since $b \in [0, 1]$.
\item $(k, l) = (2, 0)$ or $(k, l) = (0, 2)$. Then $\omega ^k + \omega^l = 1 + \omega^2 = -\omega$ and $\omega^k \omega^l = \omega^2$, so the inequality \eqref{desiredineq} can be written as
$$b^2  - 2b \re \left (u \omega \right ) + 2 \re \left (u^2\omega^2 \right ) \leq 1.$$
Reasoning similarly as before, we get
$$b^2  - 2b \cos \left (\arccos{\frac{b}{2}} + \frac{\pi}{3} \right ) + 2 \cos \left (2\arccos{\frac{b}{2}}+\frac{2\pi}{3} \right )$$
$$=b^2 - 2b \left( \frac{b}{4} - \frac{\sqrt{3}}{2} \sqrt{1 - \frac{b^2}{4}} \right) + 2 \left( -\frac{1}{2} \left( \frac{b^2}{2} - 1 \right) - \frac{\sqrt{3}}{2} \sqrt{b^2 - \frac{b^4}{4}} \right)$$
$$=b^2 - \frac{b^2}{2} + b\sqrt{3\left ( 1 - \frac{b^2}{4} \right )} - \frac{b^2}{2} + 1 - b\sqrt{3\left ( 1 - \frac{b^2}{4} \right )} = 1,$$
so again the expression turns out to be constantly equal to $1$. 

\item $(k, l)=(1, 1)$. Then \eqref{desiredineq} rewrites as
$$b^2 + 4b\re(u\omega) + 2 \re(u^2 \omega^2) \leq 1.$$
Again, we calculate
$$b^2 + 4b\re(u\omega) + 2 \re(u^2 \omega^2) = b^2 + 4b \cos \left ( \arccos{\frac{b}{2}} + \frac{\pi}{3}  \right ) + 2 \cos \left ( 2\arccos{\frac{b}{2}} + \frac{2\pi}{3}  \right )$$
$$=b^2 + 4b \left ( \frac{b}{4} - \frac{\sqrt{3}}{2} \sqrt{1 - \frac{b^2}{4}} \right ) + 2 \left ( - \frac{1}{2}  \left ( \frac{b^2}{2} - 1 \right )  -  \frac{\sqrt{3}}{2} \sqrt{b^2 - \frac{b^4}{4}} \right )$$
$$= \frac{3}{2}b^2 - 3b\sqrt{3\left ( 1 - \frac{b^2}{4} \right )} + 1 = \frac{3}{2} b \left ( b - \sqrt{12-3b^2} \right ) + 1.$$
This is clearly at most $1$, since $b \leq 1 < 3 \leq \sqrt{12-3b^2}$.
\item $(k, l)=(2, 2)$. In this case, \eqref{desiredineq} is equivalent to
$$b^2 + 4b\re(u\omega^2) + 2 \re(u^2 \omega) \leq 1.$$
Repeating a similar calculation as in the previous cases leads to
$$b^2 + 4b\re(u\omega^2) + 2 \re(u^2 \omega) = b^2 + 4b\cos \left ( \arccos{\frac{b}{2}} + \pi  \right ) + 2 \cos \left ( 2\arccos{\frac{b}{2}} \right )$$
$$=b^2 + 4b \left ( -\frac{b}{2} \right ) + 2 \left ( \frac{b^2}{2}-1 \right )=-2,$$
so in this situation we end up with a function constantly equal to $-2$.

\end{enumerate}

We have analyzed all possible cases for the pairs $(k, l)$, and therefore, the proof is finished.
\end{proof}

\end{document}
